\origpage{90}{ROBERTSON AND SEYMOUR}
\noindent $d(s, t)$ is the minimum of the number of points of $F$ in the drawing, taken over all subsets $F$ in the plane homeomorphic to $[0, 1]$ with ends $s$ and $t$. We shall speak of the distance function of $H$ without reference to any drawing, but we assume implicitly that there is given some fixed drawing of $H$. If $u, v \in V(H)$, we define $d(u, v) = d(s, t)$, where $s, t$ are the points of the drawing representing $u, v$. If $H$ and $H'$ are both walls and $H'$ is a subgraph of $H$, let the corresponding distance functions (using the drawing of $H'$ obtained from the given drawing of $H$) be $d, d'$. If $d'(s, t) \ge d(s, t) - k$ for all points $s, t$ of the plane, we write $d' \ge d - k$.

Given a wall $H$ of height $h$, we can obtain another wall of height $h - 2$ as indicated in Fig. 2, omitting half the horizontal edges of two consecutive rows and all the vertical edges in four consecutive columns and the four single-edge components that result. By generalizing this construction, we deduce

\stmt
(7.3) \textit{Let $H$ be a wall of height $h$ with distance function $d$ and let $v \in V(H)$. Let $r \ge 0$ be an integer with $h \ge 2r + 4$. Then there is a subgraph $H'$ of $H$, which is a wall of height $h - 2r - 2$, such that}

\textit{(i) $d' \ge d - 4r - 4$, where $d'$ is the distance function in $H'$}

\textit{(ii) there is no vertex $u$ of $H'$ with $d(u, v) \le r$.}

\del{The proof is left to the reader.} (The most difficult case is when $v$ is close to the infinite region of $H$, when it is necessary to delete a band of $H$ around the outside.)

\fix{We prove (7.3). Let $k = r + 1$, so that $h \ge 2k + 2$. We need a subgraph $H'$ of $H$ which is a wall of height $h - 2k$, such that $d' \ge d - 4k$ and $d(u, v) \ge k$ for every vertex $u$ of $H'$. The distances $d$ and $d'$ depend only on the drawings of $H$ and $H'$ as subsets of the plane. A homeomorphism of the plane maps subsets homeomorphic to $[0, 1]$ onto such subsets, and so applying it to the drawing of $H$ changes neither the inequality $d' \ge d - 4k$ nor the numbers $d(u, v)$. Hence we may assume that $H$ is drawn exactly as in Fig. 1, in the coordinates below.

We proceed as follows. In (1) and (2) we express $d$ through the regions of $H$. In (6) we give a test which shows $d' \ge d - 2c$ for a subwall $H'$, by comparing the regions of $H'$ with those of $H$. Then we build three kinds of subwalls and check the test for each: a copy of a smaller wall moved inside $H$, used when $v$ is near a side of $H$, a wall with a small hole, used when $r = 0$, and a wall with a large hole around $v$, used when $r \ge 1$ and $v$ is far from the sides.

The elementary wall of height $h$ has as vertices the integer points $(x, y)$ with $0 \le x \le 2h + 1$ and $0 \le y \le h$, other than $(0, 0)$ and $(2h + 1, h)$. Its \textit{horizontal} edges join $(x, y)$ and $(x + 1, y)$, and its \textit{vertical} edges join $(x, y)$ and $(x, y + 1)$ when $x + y$ is odd. Each edge is drawn as the straight segment between its ends, and this is the drawing of Fig. 1. We say that $(x, y)$ is in \textit{row} $y$ and \textit{column} $x$, and that the horizontal edge from $(x, y)$ to $(x + 1, y)$ and the vertical edge from $(x, y)$ to $(x, y + 1)$ are \textit{at} $(x, y)$. The vertices and the edges of the elementary wall are its \textit{elements}, where an edge, as an element, is the open segment without its ends. The \textit{midpoint} of an element is the vertex itself or the middle of the segment. The wall $H$ is drawn as the elementary wall, with the vertices of the subdivision placed on the segments, and so every point of the drawing of $H$ lies in exactly one element. Each wall $H' \subseteq H$ that we construct has its branch vertices at vertices of the elementary wall, and its drawing is a union of elements. We say that $H'$ \textit{keeps} the elements in its drawing and \textit{deletes} the others.

For integers $x, y$ the \textit{square} $(x, y)$ is the open unit square $(x, x + 1) \times (y, y + 1)$, and the \textit{strip} $y$ is the set of points whose second coordinate lies strictly between $y$ and $y + 1$. A \textit{position} is a pair $(X, Y)$ of integers with $X + Y$ even. The position of the square $(x, y)$ is $(x, y)$ if $x + y$ is even, and $(x + 1, y)$ if $x + y$ is odd. A position $(X, Y)$ is a \textit{brick of height} $h$ if $0 \le Y \le h - 1$ and $1 \le X \le 2h$. In this case there is no vertical edge at $(X, Y)$, since $X + Y$ is even, and so the open rectangle $(X - 1, X + 1) \times (Y, Y + 1)$ is a region of $H$. These regions are the bricks of Fig. 1, and we call this region the brick $(X, Y)$ as well. It consists of the two squares with position $(X, Y)$, namely $(X - 1, Y)$ and $(X, Y)$, and the open segment between them. Hence the regions of $H$ are the bricks of height $h$ and the infinite region, and each square lies in the brick of its position if this position is a brick of height $h$, and in the infinite region otherwise. The perimeter of $H$ consists of the points of the drawing with first coordinate at most $1$ or at least $2h$, or with second coordinate $0$ or $h$.

A region is \textit{at} a point $s$ if $s$ is in its closure. The regions at a point of an element are the regions containing the squares whose closures contain the midpoint of the element. We call these squares the squares \textit{at} the element, and these regions the regions at the element. If the midpoint is $(p, q)$, the squares at the element are the squares $(u, w)$ with $\lceil p \rceil - 1 \le u \le \lfloor p \rfloor$ and $\lceil q \rceil - 1 \le w \le \lfloor q \rfloor$. So their positions $(X, Y)$ satisfy $|X - p| \le 1$ and $q - 1 \le Y \le q$. Moreover, some square at the element has a position with $X \le p - \frac{1}{2}$, some has $X \ge p + \frac{1}{2}$, some has $Y \le q - \frac{1}{2}$ and some has $Y \ge q - \frac{1}{2}$. For the last two take $w = \lceil q \rceil - 1$ and $w = \lfloor q \rfloor$. For the first two, at a vertex $(x, y)$ one of the squares $(x - 1, y - 1)$ and $(x - 1, y)$ has position $x - 1$, and one of $(x, y - 1)$ and $(x, y)$ has position $x + 1$. At the horizontal edge at $(x, y)$ the squares $(x, y - 1)$ and $(x, y)$ have positions $x$ and $x + 1$ in some order, and at the vertical edge at $(x, y)$ the squares $(x - 1, y)$ and $(x, y)$ have positions $x - 1$ and $x + 1$. For every element at $(x, y)$, the square $(x, y)$ is at it.

For a pair $w = (w_1, w_2)$ of integers let $\|w\| = \max(|w_2|, \frac{1}{2}(|w_1| + |w_2|))$. Each of the two terms of this maximum is unchanged when $w$ is replaced by $-w$, and its value at $w + w'$ is at most the sum of its values at $w$ and at $w'$. So $\|{-w}\| = \|w\|$ and $\|w + w'\| \le \|w\| + \|w'\|$. If $w$ is the difference of two positions then $w_1 + w_2$ is even, and so $\|w\|$ is an integer. For a position $F = (X, Y)$ let
\begin{equation*}
\sigma_h(F) = \min\left(Y + 1,\ h - Y,\ \lfloor (X + 1)/2 \rfloor,\ \lfloor (2h + 2 - X)/2 \rfloor\right).
\end{equation*}
We call the four numbers in this minimum the \textit{bottom}, \textit{top}, \textit{left} and \textit{right terms} of $F$, for height $h$. By (2) below, $\|F - G\|$ is the number of steps between touching bricks from $F$ to $G$, and $\sigma_h(F)$ is the number of steps from $F$ to the infinite region. If $F$ is a brick of height $h$, each term is at least $1$, and $\sigma_h(F) = 1$ exactly when $Y \in \{0, h - 1\}$ or $X \in \{1, 2, 2h - 1, 2h\}$, that is, exactly when the closure of $F$ meets the perimeter. If $F$ is not a brick of height $h$, then one of its terms is at most $0$, and so $\sigma_h(F) \le 0$. We use two simple facts. First, for integers $A \ge B$ we have $\lfloor A/2 \rfloor - \lfloor B/2 \rfloor \le \lceil (A - B)/2 \rceil$, since $\lfloor A/2 \rfloor \le A/2$ and $\lfloor B/2 \rfloor \ge (B - 1)/2$. Second, to show that $\sigma_h(F) \le \sigma_{h'}(F') + c$ for positions $F, F'$ and heights $h, h'$, it is enough to show that each term of $F'$ for height $h'$, plus $c$, is at least some term of $F$ for height $h$.

Let $J$ be a graph drawn in the plane. Two regions of $J$ \textit{touch} if their closures have a common point. For regions $R, R'$ let $D_J(R, R')$ be the least $n \ge 0$ such that there are regions $R = R_0, R_1, \dots, R_n = R'$ with $R_{i-1}$ touching $R_i$ for $1 \le i \le n$. Thus $D_J$ satisfies the triangle inequality. We define $d_J$ from the drawing of $J$ in the same way as $d$ is defined from the drawing of $H$. The \textit{depth} of a region $R$ is $D_J(R, R_\infty)$, where $R_\infty$ is the infinite region. For a point $s$ let $[s] = 1$ if $s$ is a point of the drawing, and $[s] = 0$ otherwise.

\hspace*{12pt}(1) \textit{For distinct points $s, t$ we have $d_J(s, t) = [s] + [t] + \min D_J(R, R')$, where the minimum is over all regions $R$ at $s$ and $R'$ at $t$.}

Let $F$ be as in the definition of $d_J$, with the fewest points in the drawing. Since there is such a set with finitely many points in the drawing, $F$ has finitely many. Between two consecutive ones, $F$ runs inside one region, and the first and last of these regions are at $s$ and at $t$. At each point of the drawing other than $s$ and $t$, $F$ passes from one region into a region that touches it (or into the same region). So the regions met by $F$, in order, give a sequence of touching regions from a region at $s$ to a region at $t$, with at most as many steps as $F$ has points in the drawing other than $s$ and $t$. This proves ``$\ge$''. Conversely, choose a region $R$ at $s$, a region $R'$ at $t$, and regions $R = R_0, \dots, R_n = R'$ as above, with $n$ as small as possible, and for $1 \le i \le n$ choose a point $x_i$ of the drawing in the closures of $R_{i-1}$ and $R_i$. The regions $R_i$ are distinct, and the points $s, x_1, \dots, x_n, t$ are distinct, since otherwise the sequence could be shortened (if $x_i = x_j$ with $i < j$, then $R_{i-1}$ touches $R_j$, and if $x_1 = s$ then $R_1$ is at $s$). Join each two consecutive points of $s, x_1, \dots, x_n, t$ by an arc inside the corresponding region $R_i$, apart from its ends. The union of these arcs is a set $F$ as in the definition, and it meets the drawing in $[s] + [t] + n$ points. This proves ``$\le$'', and so proves (1).

\hspace*{12pt}(2) \textit{Two distinct bricks $F, G$ of $H$ touch if and only if $\|F - G\| = 1$, and a brick $F$ touches the infinite region if and only if $\sigma_h(F) = 1$. The depth of a brick $F$ is $\sigma_h(F)$, and $D_H(F, G) = \min(\|F - G\|, \sigma_h(F) + \sigma_h(G))$ for all bricks $F, G$.}

Let $F = (X, Y)$ and $G = (X', Y')$. Their closures are the rectangles $[X - 1, X + 1] \times [Y, Y + 1]$ and $[X' - 1, X' + 1] \times [Y', Y' + 1]$, which meet if and only if $|X - X'| \le 2$ and $|Y - Y'| \le 1$. Since $X - X'$ and $Y - Y'$ have the same parity and $F \ne G$, this holds exactly when $(|X - X'|, |Y - Y'|)$ is $(2, 0)$ or $(1, 1)$, that is, when $\|F - G\| = 1$. The closure of the infinite region is the infinite region together with the perimeter, and so $F$ touches the infinite region exactly when the closure of $F$ meets the perimeter, that is, when $\sigma_h(F) = 1$.

Let $F = R_0, R_1, \dots, R_n = R_\infty$ be touching regions with $n$ equal to the depth of $F$. Then $R_1, \dots, R_{n-1}$ are bricks. Between touching bricks each term changes by at most $1$, since the second coordinate changes by at most $1$ and the first by at most $2$, and so $\sigma_h$ changes by at most $1$. Since $R_{n-1}$ touches the infinite region, $\sigma_h(R_{n-1}) = 1$, and so $\sigma_h(F) \le n$. Conversely, a brick $(X, Y)$ with $Y \ge 1$ touches the brick $(X - 1, Y - 1)$ or the brick $(X + 1, Y - 1)$, one of which is a brick of height $h$ since $1 \le X \le 2h$. So $Y$ such steps lead from $F$ to a brick in row $0$, which touches the infinite region, and the depth of $F$ is at most $Y + 1$. In the same way, going up, the depth is at most $h - Y$. Going left or right by steps that change the first coordinate by $2$, until it is in $\{1, 2\}$ or in $\{2h - 1, 2h\}$, shows that the depth is at most $\lfloor (X + 1)/2 \rfloor$ and at most $\lfloor (2h + 2 - X)/2 \rfloor$. Hence the depth of $F$ is $\sigma_h(F)$.

Now let $F = R_0, \dots, R_n = G$ be touching regions with $n = D_H(F, G)$. If one of them is the infinite region, then $n \ge \sigma_h(F) + \sigma_h(G)$, by the depths. Otherwise they are all bricks, and $n \ge \|F - G\|$ by the triangle inequality, since $\|R_i - R_{i-1}\| = 1$. Conversely, $D_H(F, G) \le \sigma_h(F) + \sigma_h(G)$, by passing through the infinite region. Moreover $D_H(F, G) \le \|F - G\|$. For let $w = G - F$. If $|w_1| \ge |w_2|$, take $|w_2|$ steps that add $(\pm 1, \pm 1)$ and then $\frac{1}{2}(|w_1| - |w_2|)$ steps that add $(\pm 2, 0)$, with all signs towards $G$. If $|w_1| < |w_2|$, take $|w_2|$ steps that add $(\pm 1, \pm 1)$, all with the second sign towards $G$. The first $|w_1|$ of them have the first sign towards $G$, and the remaining ones (an even number) alternately move the first coordinate away from $X'$ and back, to the side of $X'$ on which it stays between $1$ and $2h$. In both cases all the positions passed are bricks, since their coordinates stay in the ranges of bricks, and the number of steps is $\|F - G\|$. This proves (2).

We shall use (2) in the following form. If $R$ and $R'$ are the regions of $H$ that contain squares with positions $F$ and $G$, then $R$ has depth $\max(\sigma_h(F), 0)$, and
\begin{equation*}
D_H(R, R') \ge \min(\|F - G\|, \text{depth}(R) + \text{depth}(R')).
\end{equation*}
If $R$ and $R'$ are bricks this is (2). If one of them is the infinite region, then $D_H(R, R')$ is the sum of the depths, and the infinite region has depth $0$ while $\sigma_h(F) \le 0$ if $F$ is not a brick.

We use the following form of a theorem of Whitney (H. Whitney, Congruent graphs and the connectivity of graphs, \textit{Amer. J. Math.} \textbf{54} (1932), 150--168): if two drawings in a sphere are subdivisions of the same simple 3-connected graph, then some homeomorphism of the sphere maps the first onto the second, each branch vertex to the corresponding branch vertex, and each subdivided edge onto the corresponding subdivided edge. For a wall $W$ of height $h'$, the \textit{side vertices} of $W$ are its branch vertices that correspond to the vertices $(x, y)$ of the elementary wall with $y \in \{0, h'\}$ or $x \in \{0, 2h' + 1\}$.

\hspace*{12pt}(3) \textit{Let $W$ be a wall of height $h' \ge 2$ drawn in the plane, and suppose that for each side vertex $w$ there is a half-line starting at $w$ that meets the drawing of $W$ only in $w$, such that these half-lines are mutually disjoint. Then some homeomorphism of the plane maps the drawing of the elementary wall of height $h'$ as in Fig. 1 onto the drawing of $W$, each vertex to the corresponding branch vertex, and each edge onto the corresponding path.}

Let $W^+$ be the graph obtained from the elementary wall of height $h'$ by adding a vertex $z$ adjacent to all side vertices. Add a point $\infty$ to the plane, to obtain a sphere. With $z$ drawn at $\infty$ and its edges drawn along the given half-lines, the drawing of $W$ becomes a drawing in the sphere of a subdivision of $W^+$. The drawing of Fig. 1 gives another one, with half-lines going down from the vertices in row $0$, up from the vertices in row $h'$, left from the vertices $(0, y)$ and right from the vertices $(2h' + 1, y)$ with $0 < y < h'$. These half-lines meet the drawing only in their starts, and they are mutually disjoint. We show below that $W^+$ is 3-connected. It is simple, and so by Whitney's theorem some homeomorphism of the sphere maps the second drawing onto the first, with each branch vertex and each subdivided edge going to the corresponding one. It maps $\infty$ to $\infty$, since $\infty$ is the point of $z$ in both drawings, and so it restricts to the required homeomorphism of the plane.

It remains to show that $W^+ - X$ is connected for every set $X$ of at most two vertices of $W^+$. If $z \in X$, then $W^+ - X$ is the elementary wall minus at most one vertex, which is connected since the elementary wall is 2-connected. (It is the union of the circuits around its bricks, and these can be ordered so that each shares an edge with an earlier one, row by row from the bottom and from left to right in each row. Adding a circuit that shares at least two vertices with a 2-connected graph gives a 2-connected graph.) If $z \notin X$, it is enough to show that every vertex $w \notin X$ of the elementary wall is joined by a path of $W^+ - X$ to a side vertex, since the side vertices are adjacent to $z$. If $w$ is a side vertex this holds trivially. Otherwise $w = (x, y)$ with $0 < x < 2h' + 1$ and $0 < y < h'$, and three paths join $w$ to side vertices and meet only in $w$: the path along row $y$ to $(0, y)$, the path along row $y$ to $(2h' + 1, y)$, and a staircase. If $x + y$ is odd the staircase is $(x, y), (x, y + 1), (x + 1, y + 1), (x + 1, y + 2), \dots$, and if $x + y$ is even it is $(x, y), (x, y - 1), (x - 1, y - 1), (x - 1, y - 2), \dots$, in both cases stopped at the first side vertex. These are paths of the wall, since each vertical step starts at a vertex with odd coordinate sum (going up) or ends at one (going down), and before the first side vertex all the vertices passed are inside the wall. The staircase leaves row $y$ at its first step. At least one of the three paths avoids $X$, since $|X| \le 2$ and $w \notin X$. This proves (3).

If (3) applies to a wall $H' \subseteq H$ of height $h'$, we name the regions of $H'$ by the regions of the elementary wall of height $h'$ drawn as in Fig. 1, using the homeomorphism. Thus the regions of $H'$ are named by the bricks of height $h'$ and by the infinite region. The region named by the brick $B'$ is the inside of the circuit $C(B')$ of $H'$ that corresponds to the circuit around $B'$, and the infinite region of $H'$ is named by the infinite region. A homeomorphism preserves touching, and so (2) for height $h'$ gives $D_{H'}$ and the depths in $H'$. Moreover every point of the drawing of $H'$ is at two different regions of $H'$, since this is true for the drawing of Fig. 1. Each region $F$ of $H$ is connected and disjoint from the drawing of $H'$, and so it lies in a region $\pi(F)$ of $H'$. In particular $\pi(R_\infty)$ is the infinite region of $H'$, since $R_\infty$ is unbounded.

All our walls $H' \subseteq H$ of height $h'$ have the following form. There are integers $R(0) < R(1) < \dots < R(h')$ such that the branch vertex of $H'$ corresponding to $(x, y)$ is a vertex of $H$ in row $R(y)$, whose column increases with $x$. Each horizontal edge of $H'$ is the path of $H$ along a row between its ends. The vertical edges of $H'$ are mutually disjoint paths of $H$. The vertical edge of $H'$ at $(x, y)$ meets rows $R(y)$ and $R(y + 1)$ of $H$ only in its ends, and it contains exactly one vertical edge of $H$ in each strip $t$ with $R(y) \le t < R(y + 1)$, and no other vertical edges of $H$. We write $c_t(x)$ for the column of this vertical edge of $H$ in strip $t$. Since the rows $R(y)$ are distinct, the columns of the branch vertices increase along each row, and each vertical edge of $H'$ lies between the heights of the two rows it joins and meets these rows only in its ends, these paths of $H$ meet only in common ends, and no branch vertex is an inner vertex of one of them. So $H'$ is a wall of height $h'$. Moreover (3) applies to $H'$, with half-lines going down from its branch vertices in row $0$, up from those in row $h'$, left from the first and right from the last branch vertex of every other row. For the drawing of $H'$ lies between the horizontal lines at heights $R(0)$ and $R(h')$, and its points at height $R(y)$ are the points of the path of $H$ along row $R(y)$ between the first and the last branch vertex of row $y$.

Finally, in each of our walls there are an integer $c_0$, and for every strip $t$ integers $y_t$ and $a_t \le a'_t$ with $a_t \equiv a'_t \equiv t - y_t \pmod 2$, such that the following holds. If $t < R(0)$ then $y_t < 0$, and if $t \ge R(h')$ then $y_t \ge h'$. If $R(0) \le t < R(h')$, then $R(y_t) \le t < R(y_t + 1)$, and $c_t(x) = x + a_t$ if $x \le c_0 - 1$ and $c_t(x) = x + a'_t$ if $x \ge c_0$, for every $x$ such that $H'$ has a vertical edge at $(x, y_t)$.

\hspace*{12pt}(4) \textit{Let $H'$ be a wall of this form, and let $F = (X, Y)$ be a brick of $H$. Let $z = X - a_Y$ if $X - a_Y \le c_0 - 1$, let $z = X - a'_Y$ if $X - a'_Y \ge c_0$, and let $z = c_0 - 1$ otherwise. Let $P(F)$ be the position of the square $(z, y_Y)$. Then $\pi(F)$ is the region of $H'$ named by $P(F)$ if $P(F)$ is a brick of height $h'$, and the infinite region of $H'$ otherwise. Moreover $F - P(F) = (\Delta X, Y - y_Y)$, where $\Delta X = a_Y$ in the first case, $\Delta X = a'_Y$ in the second case, and $\Delta X \in \{X - c_0, X - c_0 + 1\}$ in the third case. In all cases $a_Y \le \Delta X \le a'_Y$.}

The first two cases cannot both hold, since $a_Y \le a'_Y$. Let $s$ be the square $(X, Y)$, so that $s \subseteq F$ and $\pi(F)$ is the region of $H'$ that contains $s$. Let $y = y_Y$, and let $x$ be the integer with $x + y$ odd and $x + 1 \in \{z, z + 1\}$. Thus $P(F) = (x + 1, y)$. We first prove the last two statements. Since $X + Y$ and $x + 1 + y$ are even, $\Delta X = X - x - 1 \equiv Y - y \equiv a_Y \equiv a'_Y \pmod 2$. In the first case $X - z = a_Y$ and $\Delta X \in \{a_Y, a_Y - 1\}$, so $\Delta X = a_Y$ by parity. In the same way $\Delta X = a'_Y$ in the second case. In the third case $\Delta X \in \{X - c_0, X - c_0 + 1\}$ since $z = c_0 - 1$, and $X - a_Y \ge c_0$ and $X - a'_Y \le c_0 - 1$, and so $a_Y \le \Delta X \le a'_Y$.

If $Y < R(0)$ or $Y \ge R(h')$, then the horizontal line through the centre of $s$ is disjoint from the drawing of $H'$. This line and $s$ form an unbounded connected set disjoint from the drawing of $H'$, and so $s$ lies in the infinite region of $H'$. Moreover $y < 0$ or $y \ge h'$, and so $P(F)$ is not a brick of height $h'$.

Now let $R(0) \le Y < R(h')$. Then $R(y) \le Y < R(y + 1)$, and so $0 \le y \le h' - 1$. Let $\ell$ be the half-line going left from the centre $(X + \frac{1}{2}, Y + \frac{1}{2})$ of $s$. For a brick $B' = (x' + 1, y')$ of height $h'$, the circuit $C(B')$ consists of the vertical edges of $H'$ at $(x', y')$ and $(x' + 2, y')$ and of parts of rows $R(y')$ and $R(y' + 1)$ of $H$. The half-line $\ell$ meets $C(B')$ only in vertical edges of $H$ in strip $Y$, and it crosses each of them whose column is at most $X$. If $y' \ne y$, then $C(B')$ contains no vertical edge of $H$ in strip $Y$. If $y' = y$, then the vertical edges of $H$ in strip $Y$ in $C(B')$ are in columns $c_Y(x') < c_Y(x' + 2)$. A point not on a closed polygon is inside it if and only if a half-line from the point that meets the polygon only in crossings of its segments crosses it an odd number of times. Hence $s$ lies in the region named by $(x' + 1, y)$ if $c_Y(x') \le X < c_Y(x' + 2)$, and if there is no such $x'$, then $s$ lies in the infinite region of $H'$.

For every integer $w$ let $g(w) = w + a_Y$ if $w \le c_0 - 1$, and $g(w) = w + a'_Y$ if $w \ge c_0$. Then $g$ is strictly increasing, since $a_Y \le a'_Y$, and $c_Y(w) = g(w)$ whenever $H'$ has a vertical edge at $(w, y)$, that is, when $0 \le w \le 2h' + 1$ and $w + y$ is odd. We claim that $g(x) \le X < g(x + 2)$. In the first case $x \le z \le c_0 - 1$, so $g(x) = x + a_Y \le z + a_Y = X$, and $g(x + 2) \ge x + 2 + a_Y > z + a_Y = X$. In the second case $x + 2 > z \ge c_0$, so $g(x + 2) = x + 2 + a'_Y > z + a'_Y = X$, and $g(x) \le x + a'_Y \le z + a'_Y = X$. In the third case $c_0 - 2 \le x \le c_0 - 1$, so $g(x) = x + a_Y < c_0 + a_Y \le X$ and $g(x + 2) = x + 2 + a'_Y \ge c_0 + a'_Y > X$. This proves the claim.

If $x \ge 0$ and $x + 2 \le 2h' + 1$, then $c_Y(x) = g(x)$ and $c_Y(x + 2) = g(x + 2)$, since $x + y$ is odd. So $s$ lies in the region named by $(x + 1, y) = P(F)$, which is a brick of height $h'$. If $x < 0$, then every $w$ such that $H'$ has a vertical edge at $(w, y)$ satisfies $w \ge x + 2$, since $w \ge 0$ and $w \equiv x \pmod 2$. So $c_Y(w) = g(w) \ge g(x + 2) > X$, and $s$ lies in the infinite region of $H'$. In this case $x + 1 \le 0$, and $P(F)$ is not a brick of height $h'$. If $x + 2 > 2h' + 1$, then every such $w$ satisfies $w \le x$, since $w \le 2h' + 1 \le x + 1$ and $w \equiv x \pmod 2$. So $c_Y(w) = g(w) \le g(x) \le X$, and again $s$ lies in the infinite region. In this case $x + 1 \ge 2h' + 1$, and $P(F)$ is not a brick of height $h'$. This proves (4).

For a brick $F$ of $H$ we call $F - P(F)$ the \textit{move} of $F$. By (4), for each brick $F$ the depth of $\pi(F)$ is at least $\sigma_{h'}(P(F))$. For if $P(F)$ is a brick of height $h'$, then $\pi(F)$ is named by it and (2) for height $h'$ applies, and otherwise $\sigma_{h'}(P(F)) \le 0$.

\hspace*{12pt}(5) \textit{Let $F$ be a brick of $H$ whose move $(\Delta X, \Delta Y)$ satisfies $0 \le \Delta Y \le h - h'$ and $0 \le \Delta X \le 2(h - h')$. Then the bottom, top, left and right terms of $F$ for height $h$ exceed those of $P(F)$ for height $h'$ by at most $\Delta Y$, $h - h' - \Delta Y$, $\lceil \Delta X/2 \rceil$ and $\lceil (2(h - h') - \Delta X)/2 \rceil$, respectively.}

Let $F = (X, Y)$, so $P(F) = (X - \Delta X, Y - \Delta Y)$. The bottom terms are $Y + 1$ and $Y - \Delta Y + 1$, and the top terms are $h - Y$ and $h' - Y + \Delta Y$, which gives the first two bounds. The left terms are $\lfloor A/2 \rfloor$ and $\lfloor B/2 \rfloor$ with $A = X + 1$ and $B = A - \Delta X$, and the right terms are $\lfloor A/2 \rfloor$ and $\lfloor B/2 \rfloor$ with $A = 2h + 2 - X$ and $B = 2h' + 2 - X + \Delta X = A - (2(h - h') - \Delta X)$. In both cases $A \ge B$, and the first simple fact gives the last two bounds. This proves (5).

We call the four bounds in (5) the \textit{drops} of the four terms. The following is the test mentioned above.

\hspace*{12pt}(6) \textit{Let $H' \subseteq H$ be a wall of height $h'$ to which (3) applies, and let $P$ map the bricks of $H$ to positions, such that $\pi(F)$ is the region named by $P(F)$ if $P(F)$ is a brick of height $h'$, and the infinite region otherwise. Let $c \ge 1$ be an integer such that}

\hspace*{12pt}(a) \textit{$\|F - G\| \le \|P(F) - P(G)\| + 2c - 2$ for all bricks $F, G$ of $H$,}

\hspace*{12pt}(b) \textit{$\sigma_h(F) \le \sigma_{h'}(P(F)) + c$ for every brick $F$ of $H$, and}

\hspace*{12pt}(c) \textit{every element deleted by $H'$ is at a region $F$ of $H$ whose depth is at most the depth of $\pi(F)$ plus $c - 1$.}

\textit{Then $d' \ge d - 2c$.}

As above, the depth of $\pi(F)$ is at least $\sigma_{h'}(P(F))$ for every brick $F$. So by (b) every region $F$ of $H$ has depth at most the depth of $\pi(F)$ plus $c$. For the infinite region this holds since both depths are $0$. Let $s \ne t$ be points (for $s = t$ both distances are $0$), write $[s]'$ for $[s]$ taken in $H'$, and choose regions $F', G'$ of $H'$ at $s, t$ such that $d'(s, t) = [s]' + [t]' + D_{H'}(F', G')$, by (1) for $H'$. Since $F'$ is open and the drawing of $H$ is a finite union of segments, $F'$ contains points arbitrarily close to $s$ that are not in the drawing of $H$. There are finitely many regions of $H$, and so one region $F$ of $H$ contains such points arbitrarily close to $s$. Then $F$ is at $s$, and $F \subseteq F'$ since $F$ is connected and disjoint from the drawing of $H'$. If $s$ is a point of $H$ but not of $H'$, then a small disc around $s$ meets no point of $H'$, since the drawing of $H'$ is closed. This disc lies in $F'$, and so all the regions of $H$ at $s$ lie in $F'$. In this case $s$ lies in an element deleted by $H'$, and by (c) we may choose $F$ with depth at most that of $F'$ plus $c - 1$. In every case $[s] - [s]' \le 1$, and $[s] - [s]'$ plus the depth of $F$ is at most the depth of $F'$ plus $c$. If $s$ is a point of $H$ but not of $H'$ this is the choice just made, and otherwise $[s] = [s]'$ and it follows from the second sentence of this proof. Choose $G$ in the same way for $t$. By (1) for $H$, $d(s, t) \le [s] + [t] + D_H(F, G)$. By (2) for $H'$, either $D_{H'}(F', G')$ is the sum of the depths of $F'$ and $G'$, or $F', G'$ are named by bricks $B_1, B_2$ of height $h'$ and $D_{H'}(F', G') = \|B_1 - B_2\|$. In the first case,
\begin{align*}
d(s, t) &\le [s] + [t] + \text{depth}(F) + \text{depth}(G)\\
&\le [s]' + [t]' + \text{depth}(F') + \text{depth}(G') + 2c = d'(s, t) + 2c.
\end{align*}
In the second case $F, G$ are bricks, since $\pi$ maps the infinite region of $H$ to that of $H'$. Moreover $\pi(F) = F'$ and $\pi(G) = G'$, so $P(F) = B_1$ and $P(G) = B_2$. By (2) for $H$ and by (a),
\begin{align*}
d(s, t) &\le [s] + [t] + \|F - G\|\\
&\le ([s]' + 1) + ([t]' + 1) + \|B_1 - B_2\| + 2c - 2 = d'(s, t) + 2c.
\end{align*}
This proves (6).

By (4), the map $P$ defined there satisfies the hypothesis of (6) for every wall of our form. We check (a), (b), (c) as follows. If all the moves of bricks of $H$ satisfy the hypothesis of (5), and all the drops are at most $c$, then (b) holds, by the second simple fact. For (c) and an element $m$ deleted by $H'$, it is enough to find a square at $m$ whose position $F$ is not a brick of height $h$, or satisfies $\sigma_h(F) \le c - 1$, or satisfies $\sigma_h(F) \le \sigma_{h'}(P(F)) + c - 1$. This is because the region of $H$ containing the square lies at $m$. It is the infinite region if $F$ is not a brick of height $h$, and otherwise it is the brick $F$, of depth $\sigma_h(F)$, while $\pi(F)$ has depth at least $\max(\sigma_{h'}(P(F)), 0)$. For the last inequality it is enough that each term of $P(F)$, plus $c - 1$, is at least some term of $F$.

For (ii) we use (1) as follows. Let $m_v$ be the element of $H$ that contains $v$. The regions of $H$ at $v$ are the regions at $m_v$, and similarly for every point of the drawing. Let $u$ be a vertex of $H'$ with $u \ne v$, such that every region of $H$ at $u$ has distance $D_H$ at least $k - 2$ from every region of $H$ at $v$. Then $d(u, v) \ge 1 + 1 + (k - 2) = k$ by (1). If $H'$ deletes $m_v$, then $u \ne v$ for every vertex $u$ of $H'$, since $u$ lies in an element kept by $H'$.

We use the rotation $\rho$ of the plane given by $\rho(x, y) = (2h + 1 - x, h - y)$. Since $h$ is even, $\rho$ maps each element of $H$ onto an element: the vertex $(x, y)$ to the vertex $(2h + 1 - x, h - y)$, the horizontal edge at $(x, y)$ to the horizontal edge at $(2h - x, h - y)$, and the vertical edge at $(x, y)$ to the vertical edge at $(2h + 1 - x, h - 1 - y)$, which exists because $(2h + 1 - x) + (h - 1 - y) = 3h - x - y$ is odd. It maps the square $(x, y)$ to the square $(2h - x, h - 1 - y)$, and the brick $(X, Y)$ to the brick $(2h + 1 - X, h - 1 - Y)$. We let $\rho$ act on all positions by this formula. It exchanges the bottom and top terms of positions, and the left and right terms, and so it preserves $\sigma_h$, and also $\|F - G\|$ for positions $F, G$. Since $\rho$ maps the drawing of $H$ onto itself, it maps regions of $H$ to regions of $H$ and preserves touching and depths. Let $\rho'$ be defined in the same way with $h'$ in place of $h$. Let $H' \subseteq H$ and $P$ be as in (6). Then the subgraph $\rho(H')$ of $H$, whose drawing is the image under $\rho$ of the drawing of $H'$, is a wall of height $h'$ whose branch vertex corresponding to $(x, y)$ is the image under $\rho$ of the branch vertex of $H'$ corresponding to $\rho'(x, y)$. Composing the homeomorphism of (3) for $H'$ with $\rho$ and $\rho'$ shows that (3) applies to $\rho(H')$, and that the map $F \mapsto \rho'(P(\rho(F)))$ has the property of $P$ in (6) for $\rho(H')$. Moreover $\rho(H')$ deletes an element $m$ if and only if $H'$ deletes $\rho(m)$. Hence if $H'$ satisfies (a), (b), (c) of (6) for some $c$, then so does $\rho(H')$. We call $\rho(H')$ the \textit{turned} $H'$.

We come to the subwalls. Let $h' = h - 2k$. Let $1 \le i \le 4k - 1$ and $1 \le j \le 2k - 1$ with $i + j$ even. The \textit{moved wall} $T(i, j)$ is the subgraph of $H$ whose drawing is the drawing of the elementary wall of height $h'$, as in Fig. 1, translated by $(i, j)$. Each vertex $(x, y)$ of the elementary wall of height $h'$ is translated to the vertex $(x + i, y + j)$ of $H$, since $1 \le x + i \le 2h' + 1 + 4k - 1 = 2h$ and $1 \le y + j \le h' + 2k - 1 = h - 1$. Each vertical edge is translated to a vertical edge of $H$, since $i + j$ is even. So $T(i, j)$ is a wall of our form, with $R(y) = y + j$ and $c_t(x) = x + i$, and we take $c_0 = 0$ and $a_t = a'_t = i$ and $y_t = t - j$ for every strip $t$. By (4), $P(F) = F - (i, j)$ for every brick $F$ of $H$. The wall $T(i, j)$ keeps exactly the elements whose midpoint minus $(i, j)$ is the midpoint of an element of the elementary wall of height $h'$. In particular, the midpoint $(p, q)$ of every element kept by $T(i, j)$ satisfies $i \le p \le i + 2h' + 1$ and $j \le q \le j + h'$.

\hspace*{12pt}(7) \textit{The moved wall $T(i, j)$ satisfies (a), (b), (c) of (6) with $c = 2k$, and so $d' \ge d - 4k$ for $H' = T(i, j)$.}

Since all moves equal $(i, j)$, we have $\|F - G\| = \|P(F) - P(G)\|$, which gives (a). By (5) with $h - h' = 2k$, the drops are $j$, $2k - j$, $\lceil i/2 \rceil$ and $\lceil (4k - i)/2 \rceil$, all at most $2k$. This gives (b). For (c), let $m$ be an element deleted by $T(i, j)$, with midpoint $(p, q)$. Let the \textit{infinite wall} have all integer points as vertices, with horizontal and vertical edges defined as for the elementary wall. Translating by $-(i, j)$ maps $m$ to a vertex or edge of the infinite wall, since $i + j$ is even. A vertex or edge of the infinite wall whose midpoint $(p'', q'')$ satisfies $0 \le p'' \le 2h' + 1$ and $0 < q'' < h'$ is an element of the elementary wall of height $h'$, since the vertices and edges of the infinite wall in this box that are missing from the elementary wall have $q'' \in \{0, h'\}$. Hence $q \le j$, or $q \ge j + h'$, or $p \le i - \frac{1}{2}$, or $p \ge i + 2h' + \frac{3}{2}$, since $p$ is a multiple of $\frac{1}{2}$. In the first case choose a square at $m$ whose position $F = (X, Y)$ has $Y \le q - \frac{1}{2}$, so that $Y \le j - 1$. The bottom term of $F$ is $Y + 1 \le j \le 2k - 1$, and the second coordinate of $P(F)$ is $Y - j < 0$. In the second case choose a square with $Y \ge q - \frac{1}{2}$, so that $Y \ge j + h'$. The top term of $F$ is $h - Y \le 2k - j \le 2k - 1$, and the second coordinate of $P(F)$ is at least $h'$. In the third case choose a square with $X \le p - \frac{1}{2} \le i - 1$. The left term of $F$ is at most $\lfloor i/2 \rfloor \le 2k - 1$, and the first coordinate of $P(F)$ is negative. In the fourth case choose a square with $X \ge p + \frac{1}{2} \ge i + 2h' + 2$. The right term of $F$ is at most $\lfloor (2h - 2h' - i)/2 \rfloor = \lfloor (4k - i)/2 \rfloor \le 2k - 1$, and the first coordinate of $P(F)$ is at least $2h' + 2$. In each case $\sigma_h(F) \le 2k - 1$, which gives (c). By (6), $d' \ge d - 4k$. This proves (7).

Since the drawing of the elementary wall of height $h'$ is mapped onto itself by $\rho'$, the rotation $\rho$ maps the drawing of $T(i, j)$ onto the drawing of $T(4k - i, 2k - j)$. Thus $\rho(T(i, j)) = T(4k - i, 2k - j)$.

Next, let $k = 1$, so that $h' = h - 2$. We use a variant of the construction of Fig. 2, turned upside down. Let $Y_0, c_0$ be integers with $1 \le Y_0 \le h'$ and $1 \le c_0 \le 2h' + 1$. The wall $S(Y_0, c_0)$ is the subgraph of $H$ defined as follows. Its branch vertex corresponding to $(x, y)$ is the vertex $(x, y)$ of $H$ if $y \le Y_0 - 1$ and $x \le c_0 - 1$, the vertex $(x + 4, y)$ if $y \le Y_0 - 1$ and $x \ge c_0$, and the vertex $(x + 2, y + 2)$ if $y \ge Y_0$. Its horizontal edges are the paths of $H$ along rows between their ends. Its vertical edge at $(x, y)$ is the vertical edge of $H$ at the vertex of $H$ corresponding to $(x, y)$, if $y \ne Y_0 - 1$. If $y = Y_0 - 1$, then $x \equiv Y_0 \pmod 2$, and the vertical edge at $(x, y)$ is a staircase made of five edges of the elementary wall. If $x \le c_0 - 1$, it goes from $(x, Y_0 - 1)$ up to $(x, Y_0)$, right to $(x + 1, Y_0)$, up to $(x + 1, Y_0 + 1)$, right to $(x + 2, Y_0 + 1)$, and up to $(x + 2, Y_0 + 2)$. If $x \ge c_0$, it goes from $(x + 4, Y_0 - 1)$ up to $(x + 4, Y_0)$, left to $(x + 3, Y_0)$, up to $(x + 3, Y_0 + 1)$, left to $(x + 2, Y_0 + 1)$, and up to $(x + 2, Y_0 + 2)$. These vertical steps are edges, since they start at vertices with odd coordinate sum, as $x \equiv Y_0$.

\hspace*{12pt}(8) \textit{The wall $S(Y_0, c_0)$ is a wall of height $h'$ of our form, and every move is one of $(0, 0)$, $(2, 0)$, $(4, 0)$, $(1, 1)$, $(3, 1)$, $(2, 2)$. It deletes the vertical edges at $(x, y)$ with $y \le Y_0 - 1$ and $c_0 \le x \le c_0 + 3$. If $c_0 \equiv Y_0 \pmod 2$, it also deletes the vertices $(c_0, Y_0)$, $(c_0 + 1, Y_0)$, $(c_0 + 2, Y_0)$ and $(c_0 + 1, Y_0 + 1)$, the horizontal edges at $(c_0 - 1, Y_0)$, $(c_0, Y_0)$, $(c_0, Y_0 + 1)$ and $(c_0 + 1, Y_0 + 1)$, and the vertical edge at $(c_0 + 1, Y_0)$. Every element $m$ it deletes, with midpoint $(p, q)$, is of one of the following kinds:}

\hspace*{12pt}(i) \textit{a vertical edge at $(x, y)$ with $y \le Y_0 - 1$ and $c_0 \le x \le c_0 + 3$,}

\hspace*{12pt}(ii) \textit{an element with $Y_0 \le q \le Y_0 + 1$,}

\hspace*{12pt}(iii) \textit{an element with $p < 2$, or $p > 2h - 1$, or $q = h$,}

\hspace*{12pt}(iv) \textit{if $c_0 = 1$, a vertex $(x, 0)$ or a horizontal edge at $(x, 0)$ with $1 \le x \le 4$.}

The branch vertices are vertices of $H$, since $x + 4 \le 2h' + 5 = 2h + 1$ and $y + 2 \le h' + 2 = h$, and the staircases use only the rows $Y_0, Y_0 + 1 \le h - 1$. We check the form with $R(y) = y$ for $y \le Y_0 - 1$ and $R(y) = y + 2$ for $y \ge Y_0$. The columns of the branch vertices increase along each row. Each staircase meets rows $Y_0 - 1$ and $Y_0 + 2$ only in its ends and has one vertical edge of $H$ in each of the strips $Y_0 - 1$, $Y_0$, $Y_0 + 1$. In rows $Y_0$ and $Y_0 + 1$ the staircase for $x \le c_0 - 1$ uses the columns $x, x + 1$ and $x + 1, x + 2$, and the staircase for $x \ge c_0$ uses the columns $x + 3, x + 4$ and $x + 2, x + 3$. Since staircases on the same side differ by $2$ in $x$, and since the columns used on the left are at most $c_0 + 1$ and those on the right at least $c_0 + 2$, the staircases are mutually disjoint. So $S(Y_0, c_0)$ has our form. We take $c_0$ as given, and for every strip $t$ we take $y_t = t$ and $(a_t, a'_t) = (0, 4)$ if $t \le Y_0 - 1$, $y_t = Y_0 - 1$ and $(a_t, a'_t) = (1, 3)$ if $t = Y_0$, $y_t = Y_0 - 1$ and $(a_t, a'_t) = (2, 2)$ if $t = Y_0 + 1$, and $y_t = t - 2$ and $(a_t, a'_t) = (2, 2)$ if $t \ge Y_0 + 2$. These satisfy the conditions of our form, by the definition of the vertical edges, and $R(0) = 0$ and $R(h') = h$. By (4), the moves in strips $t \le Y_0 - 1$ are $(\Delta X, 0)$ with $0 \le \Delta X \le 4$ and $\Delta X$ even, those in strip $Y_0$ are $(1, 1)$ and $(3, 1)$, and those in higher strips are $(2, 2)$.

By the form, the vertical edges of $H$ that $S(Y_0, c_0)$ keeps in strips $t \le Y_0 - 1$ are in columns $c_t(x) = x \le c_0 - 1$ or $c_t(x) = x + 4 \ge c_0 + 4$, which shows that the elements of kind (i) are deleted. Now let $c_0 \equiv Y_0 \pmod 2$. The vertices and edges of $H$ in rows $Y_0$ and $Y_0 + 1$ and in strip $Y_0$ that $S(Y_0, c_0)$ keeps are those of the staircases, since the other rows and vertical edges of $S(Y_0, c_0)$ lie in rows $y \le Y_0 - 1$ or $y \ge Y_0 + 2$ and in strips other than $Y_0$. The staircases are at the points $(x, Y_0 - 1)$ with $x \equiv Y_0 \pmod 2$, so that $x \equiv c_0 \pmod 2$, and hence $x \le c_0 - 2$ or $x \ge c_0$. The staircase for $x \le c_0 - 2$ meets these rows and this strip in the vertices $(x, Y_0)$, $(x + 1, Y_0)$, $(x + 1, Y_0 + 1)$, $(x + 2, Y_0 + 1)$, the horizontal edges at $(x, Y_0)$ and $(x + 1, Y_0 + 1)$, and the vertical edge at $(x + 1, Y_0)$. The staircase for $x \ge c_0$ meets them in the vertices $(x + 4, Y_0)$, $(x + 3, Y_0)$, $(x + 3, Y_0 + 1)$, $(x + 2, Y_0 + 1)$, the horizontal edges at $(x + 3, Y_0)$ and $(x + 2, Y_0 + 1)$, and the vertical edge at $(x + 3, Y_0)$. Hence a vertex $(x', Y_0)$ or a vertical edge at $(x', Y_0)$ that is kept has $x' \le c_0 - 1$ or $x' \ge c_0 + 3$, a horizontal edge at $(x', Y_0)$ that is kept has $x' \le c_0 - 2$ or $x' \ge c_0 + 3$, a vertex $(x', Y_0 + 1)$ that is kept has $x' \le c_0$ or $x' \ge c_0 + 2$, and a horizontal edge at $(x', Y_0 + 1)$ that is kept has $x' \le c_0 - 1$ or $x' \ge c_0 + 2$. This shows that the other listed elements are deleted.

Now let $m$ be deleted and not of kinds (i) and (ii). Suppose first that $q \le Y_0 - \frac{1}{2}$. Then $m$ is a vertex or a horizontal edge in a row $y \le Y_0 - 1$, since the vertical edges of $H$ in strips $t \le Y_0 - 1$ in other columns than $c_0, \dots, c_0 + 3$ are kept. The path of row $y$ of $S(Y_0, c_0)$ runs along row $y$ of $H$ from column $0$, or from column $1$ if $y = 0$ and $c_0 \ge 2$, or from column $5$ if $y = 0$ and $c_0 = 1$, to column $2h' + 5 = 2h + 1$. So it is all of row $y$ of $H$ unless $y = 0$ and $c_0 = 1$, and $m$ is of kind (iv). Suppose next that $q \ge Y_0 + \frac{3}{2}$. The part of the drawing of $S(Y_0, c_0)$ above height $Y_0 + 1$ is the part of the drawing of the elementary wall of height $h'$ above height $Y_0 - 1$, translated by $(2, 2)$. Indeed, the rows $y \ge Y_0$ and the vertical edges at $(x, y)$ with $y \ge Y_0$ are translated by $(2, 2)$, and so is the last step of each staircase. So $m$ translated by $-(2, 2)$ is a vertex or edge of the infinite wall which is not an element of the elementary wall of height $h'$, while $0 < q - 2 \le h'$. As in the proof of (7), it follows that $p - 2 < 0$, or $p - 2 > 2h' + 1$, or $q - 2 = h'$, and so $m$ is of kind (iii). This proves (8).

We call $S(Y_0, c_0)$ \textit{good} if $c_0 \ge 2Y_0$ when $Y_0$ is odd, and $c_0 + 2Y_0 \le 2h' + 2$ when $Y_0$ is even.

\hspace*{12pt}(9) \textit{If $S(Y_0, c_0)$ is good, then it satisfies (a), (b), (c) of (6) with $c = 2$.}

Each of the six moves $\mu$ of (8) satisfies $|\mu_1 - 2| + |\mu_2| \le 2$, and so any two moves $\mu, \mu'$ satisfy $|\mu_1 - \mu'_1| + |\mu_2 - \mu'_2| \le 4$ and $|\mu_2 - \mu'_2| \le 2$. Hence $\|\mu - \mu'\| \le 2$, and so $\|F - G\| \le \|P(F) - P(G)\| + 2$ for all bricks $F, G$, since $F - G = (P(F) - P(G)) + (\mu - \mu')$ where $\mu, \mu'$ are the moves of $F, G$. This gives (a). By (5) with $h - h' = 2$, all the drops are at most $2$, since $0 \le \Delta Y \le 2$ and $0 \le \Delta X \le 4$. This gives (b).

For (c) let $m$ be a deleted element, of one of the kinds of (8). Let $(x, y)$ be such that $m$ is the vertex $(x, y)$ or an edge at $(x, y)$. If $m$ is of kind (iii) or (iv), let $F = (X, Y)$ be the position of the square $(x, y)$, which is at $m$. If $p < 2$ then $x \le 1$ and $X \le 2$, so the left term of $F$ is at most $1$. If $p > 2h - 1$ then $x \ge 2h - 1$ and $X \ge 2h - 1$, so the right term of $F$ is at most $1$. If $q = h$, then $Y = h$ and $F$ is not a brick of height $h$. If $m$ is of kind (iv), the bottom term of $F$ is $1$. So in these cases $F$ is as required.

Let $m$ be of kind (ii). Then $m$ has a square in strip $Y_0$ at it: the square $(x, Y_0)$ if $y = Y_0$, and the square $(x, Y_0)$ below $m$ if $m$ is a vertex or horizontal edge in row $Y_0 + 1$. Let $F = (X, Y_0)$ be its position. If $F$ is not a brick of height $h$ we are done, so we may assume that it is. Its move is $(1, 1)$ or $(3, 1)$. By (4), the move is $(1, 1)$ if $X \le c_0$, it is $(3, 1)$ if $X \ge c_0 + 3$, and otherwise it is $(X - c_0, 1)$ or $(X - c_0 + 1, 1)$. So if the move is $(1, 1)$ then $X \le c_0 + 1$, and if it is $(3, 1)$ then $X \ge c_0 + 2$. By (5), the bottom and top terms drop by $1$. Suppose that $Y_0$ is odd, so that $X$ is odd. If the move is $(1, 1)$, the left term drops by $\lceil 1/2 \rceil = 1$, and the right term drops from $(2h + 1 - X)/2$ to $\lfloor (2h - 1 - X)/2 \rfloor = (2h - 1 - X)/2$, that is, by $1$. If the move is $(3, 1)$, the right term drops by $\lceil 1/2 \rceil = 1$. The left term of $P(F)$ is $\lfloor (X - 2)/2 \rfloor = (X - 3)/2$, and since $X$ is odd and $X \ge c_0 + 2 \ge 2Y_0 + 2$ because $S(Y_0, c_0)$ is good, we have $X \ge 2Y_0 + 3$. So this left term plus $1$ is at least $Y_0 + 1$, the bottom term of $F$. Now suppose that $Y_0$ is even, so that $X$ is even. If the move is $(3, 1)$, the left term drops from $X/2$ to $\lfloor (X - 2)/2 \rfloor = X/2 - 1$, and the right term by $1$. If the move is $(1, 1)$, the left term drops by $1$. The right term of $P(F)$ is $\lfloor (2h' + 3 - X)/2 \rfloor$, and since $X \le c_0 + 1 \le 2h' + 3 - 2Y_0$ because $S(Y_0, c_0)$ is good, it is at least $Y_0$. So this right term plus $1$ is at least $Y_0 + 1$, the bottom term of $F$. In all cases each term of $P(F)$ plus $1$ is at least some term of $F$, as required.

Let $m$ be the vertical edge at $(x, y)$ of kind (i), so that $y \le Y_0 - 1$ and $c_0 \le x \le c_0 + 3$. Suppose first that $Y_0$ is odd, and let $F = (X, y)$ be the position of the square $(x, y)$, so that $X = x + 1$ and $c_0 + 1 \le X \le c_0 + 4$. In the definition of $P(F)$ in (4), with $a_y = 0$ and $a'_y = 4$, the first case does not occur. So $z = c_0 - 1$ if $X \le c_0 + 3$, and $z = c_0$ if $X = c_0 + 4$. In both cases $P(F) = (X', y)$ with $X' \in \{c_0 - 1, c_0\}$, and the move is $(X - X', 0)$ with $X - X' \in \{2, 4\}$, since it is positive and even. By (5), the bottom term does not drop, and the right term drops by at most $\lceil (4 - 2)/2 \rceil = 1$. The left term of $P(F)$ is $\lfloor (X' + 1)/2 \rfloor \ge \lfloor c_0/2 \rfloor \ge Y_0 \ge y + 1$, using $c_0 \ge 2Y_0$. So it is at least the bottom term of $P(F)$, which equals that of $F$. The right term of $P(F)$ is $\lfloor (2h' + 2 - X')/2 \rfloor \le \lfloor (2h' + 3 - 2Y_0)/2 \rfloor = h' + 1 - Y_0 \le h' - y$, which is the top term of $P(F)$. So the top term of $P(F)$ plus $1$ is at least its right term plus $1$, which is at least the right term of $F$. Now suppose that $Y_0$ is even, and let $F = (X, y)$ be the position of the square $(x - 1, y)$, so that $X = x - 1$ and $c_0 - 1 \le X \le c_0 + 2$. If $X = c_0 - 1$, we are in the first case of (4) and $P(F) = F$. Otherwise we are in the third case, $P(F) = (X', y)$ with $X' \in \{c_0 - 1, c_0\}$, and the move $(X - X', 0)$ satisfies $X - X' \in \{0, 2\}$. So the bottom term does not drop, and the left term drops by at most $1$. The left term of $P(F)$ is $\lfloor (X' + 1)/2 \rfloor \le \lfloor (c_0 + 1)/2 \rfloor \le h' + 1 - Y_0 \le h' - y$, using $c_0 \le 2h' + 2 - 2Y_0$, and this is the top term of $P(F)$. The right term of $P(F)$ is $\lfloor (2h' + 2 - X')/2 \rfloor \ge \lfloor (2h' + 2 - c_0)/2 \rfloor \ge Y_0 \ge y + 1$, the bottom term of $F$. In both cases each term of $P(F)$ plus $1$ is at least some term of $F$. This proves (9).

\hspace*{12pt}(10) \textit{If $k = 1$, there is a wall $H' \subseteq H$ of height $h - 2$ which satisfies (a), (b), (c) of (6) with $c = 2$ and deletes $m_v$.}

If $T(1, 1)$ or $T(3, 1)$ deletes $m_v$, we take it as $H'$, by (7). So we may assume that both keep $m_v$. Let $(x, y)$ be such that $m_v$ is the vertex $(x, y)$ or an edge at $(x, y)$. Since $T(1, 1)$ and $T(3, 1)$ keep $m_v$, the elementary wall of height $h'$ has an element of the same kind at $(x - 1, y - 1)$ and at $(x - 3, y - 1)$. Hence $3 \le x \le 2h' + 2$ and $1 \le y \le h' + 1$. Moreover $x \le 2h' + 1$ if $m_v$ is a horizontal edge, and $y \le h'$ if $m_v$ is a vertical edge. Since $h = h' + 2$, the rotation $\rho$ maps the vertex $(x, y)$ to the vertex $(2h' + 5 - x, h' + 2 - y)$, the horizontal edge at $(x, y)$ to the horizontal edge at $(2h' + 4 - x, h' + 2 - y)$, and the vertical edge at $(x, y)$ to the vertical edge at $(2h' + 5 - x, h' + 1 - y)$.

We take as $H'$ the wall $S(Y_0, c_0)$ or the turned $S(Y_0, c_0)$, for integers $Y_0, c_0$ with $c_0 \equiv Y_0 \pmod 2$ such that $S(Y_0, c_0)$ is defined and good, that is, such that $1 \le Y_0 \le h'$ and $1 \le c_0 \le 2h' + 1$, and $c_0 \ge 2Y_0$ if $Y_0$ is odd and $c_0 + 2Y_0 \le 2h' + 2$ if $Y_0$ is even. Then $H'$ satisfies (a), (b), (c) of (6) with $c = 2$, by (9) and the remark about turned walls. By (8), $S(Y_0, c_0)$ deletes $m_v$ if $m_v$ is one of the elements listed there, and the turned $S(Y_0, c_0)$ deletes $m_v$ if $\rho(m_v)$ is one of them, since it deletes $m_v$ if and only if $S(Y_0, c_0)$ deletes $\rho(m_v)$.

For each kind of $m_v$ and each parity of $x$ and $y$, we give below two or three candidates for $H'$, in order. In every candidate $c_0 \equiv Y_0 \pmod 2$. For each candidate we name the element listed in (8) that $m_v$ is, or that $\rho(m_v)$ is for a turned candidate, and we give a condition that holds if and only if $S(Y_0, c_0)$ is defined and good. We take as $H'$ the first candidate whose condition holds, and we show that the condition of the last candidate holds if the earlier ones fail. Each condition comes from putting the values of $Y_0$ and $c_0$ into the inequalities above. In this, $1 \le c_0 \le 2h' + 1$ always follows from the bounds on $x$ and the parity of $x$, and $1 \le Y_0 \le h'$ follows from the bounds on $y$ and the parity of $y$, except that $Y_0 = h' + 2 - y$ with $y$ odd needs $y \ge 3$, $Y_0 = y$ with $y$ odd needs $y \le h'$, and $Y_0 = y + 1$ needs $y \le h' - 1$. These are part of the conditions below, or follow from them and the bounds.

Let $m_v$ be a vertex. If $x$ and $y$ are even, the candidates are $S(y - 1, x - 1)$, where $m_v$ is the vertex $(c_0 + 1, Y_0 + 1)$, with the condition $x \ge 2y - 1$; then $S(y, x - 2)$, where $m_v$ is the vertex $(c_0 + 2, Y_0)$, with the condition $x + 2y \le 2h' + 4$; and then the turned $S(h' + 2 - y, 2h' + 4 - x)$, where $\rho(m_v)$ is the vertex $(c_0 + 1, Y_0)$, with the condition $x + 2y \ge 2h' + 6$. If the second condition fails, then $x + 2y \ge 2h' + 5$, and so the third holds, since $x + 2y$ is even.

If $x$ and $y$ are odd, the candidates are the turned $S(h' + 2 - y, 2h' + 4 - x)$, where $\rho(m_v)$ is the vertex $(c_0 + 1, Y_0)$, with the condition that $y \ge 3$ and $x \le 2y$; and then $S(y, x)$, where $m_v$ is the vertex $(c_0, Y_0)$, with the condition that $y \le h'$ and $x \ge 2y$. If the first condition fails, then $y = 1$ or $x > 2y$. If $y = 1$, then $y \le h'$ and $x \ge 3 > 2y$. If $x > 2y$, then $2y < x \le 2h' + 2$, and so $y \le h'$. In both cases the second condition holds.

If $x$ is odd and $y$ is even, the candidates are the turned $S(h' + 1 - y, 2h' + 4 - x)$, where $\rho(m_v)$ is the vertex $(c_0 + 1, Y_0 + 1)$, with the condition $x \le 2y + 2$; then $S(y, x - 1)$, where $m_v$ is the vertex $(c_0 + 1, Y_0)$, with the condition $x + 2y \le 2h' + 3$; and then the turned $S(h' + 2 - y, 2h' + 3 - x)$, where $\rho(m_v)$ is the vertex $(c_0 + 2, Y_0)$, with the condition $x + 2y \ge 2h' + 5$. If the second condition fails, then $x + 2y \ge 2h' + 4$, and so the third holds, since $x + 2y$ is odd.

If $x$ is even and $y$ is odd, the candidates are $S(y, x - 1)$, where $m_v$ is the vertex $(c_0 + 1, Y_0)$, with the condition $x \ge 2y + 1$, which gives $y \le h'$ since $x \le 2h' + 2$; and then the turned $S(h' + 2 - y, 2h' + 5 - x)$, where $\rho(m_v)$ is the vertex $(c_0, Y_0)$, with the condition that $y \ge 3$ and $x \le 2y + 1$. If the first condition fails, then $x \le 2y$, and so $y \ge 3$, since $x \ge 3$ and $y$ is odd. So the second condition holds.

Let $m_v$ be a horizontal edge. If $x$ and $y$ are even, the candidates are $S(y - 1, x - 1)$, where $m_v$ is the horizontal edge at $(c_0 + 1, Y_0 + 1)$, with the condition $x \ge 2y - 1$; and then the turned $S(h' + 1 - y, 2h' + 3 - x)$, where $\rho(m_v)$ is the horizontal edge at $(c_0 + 1, Y_0 + 1)$, with the condition $x \le 2y + 1$, which holds if the first fails.

If $x$ and $y$ are odd, the candidates are $S(y, x)$, where $m_v$ is the horizontal edge at $(c_0, Y_0)$, with the condition $x \ge 2y$, which gives $y \le h'$ since $x \le 2h' + 1$; and then the turned $S(h' + 2 - y, 2h' + 4 - x)$, where $\rho(m_v)$ is the horizontal edge at $(c_0, Y_0)$, with the condition that $y \ge 3$ and $x \le 2y$. If the first condition fails, then $x < 2y$, and so $y \ge 3$, since $x \ge 3$ and $y$ is odd. So the second condition holds.

If $x$ is odd and $y$ is even, the candidates are $S(y - 1, x)$, where $m_v$ is the horizontal edge at $(c_0, Y_0 + 1)$, with the condition $x \ge 2y - 2$; and then the turned $S(h' + 1 - y, 2h' + 4 - x)$, where $\rho(m_v)$ is the horizontal edge at $(c_0, Y_0 + 1)$, with the condition $x \le 2y + 2$, which holds if the first fails.

If $x$ is even and $y$ is odd, the candidates are $S(y, x + 1)$, where $m_v$ is the horizontal edge at $(c_0 - 1, Y_0)$, with the condition $x + 1 \ge 2y$, which gives $y \le h'$ since $x \le 2h'$ as $x$ is even; and then the turned $S(h' + 2 - y, 2h' + 5 - x)$, where $\rho(m_v)$ is the horizontal edge at $(c_0 - 1, Y_0)$, with the condition that $y \ge 3$ and $x \le 2y + 1$. If the first condition fails, then $x \le 2y - 2$, and so $y \ge 3$, since $x \ge 3$. So the second condition holds.

Let $m_v$ be a vertical edge, so that $x + y$ is odd. If $y$ is even, the candidates are the turned $S(h' + 1 - y, 2h' + 4 - x)$, where $\rho(m_v)$ is the vertical edge at $(c_0 + 1, Y_0)$, with the condition $x \le 2y + 2$; and then $S(y + 1, x)$, where $m_v$ is the vertical edge at $(c_0, Y_0 - 1)$, with the condition that $y \le h' - 1$ and $x \ge 2y + 2$. If the first condition fails, then $x \ge 2y + 3$, and so $2y + 3 \le 2h' + 2$, which gives $y \le h' - 1$. So the second condition holds.

If $y$ is odd, the candidates are $S(y, x - 1)$, where $m_v$ is the vertical edge at $(c_0 + 1, Y_0)$, with the condition $x \ge 2y + 1$; and then the turned $S(h' + 2 - y, 2h' + 5 - x)$, where $\rho(m_v)$ is the vertical edge at $(c_0, Y_0 - 1)$, with the condition that $y \ge 3$ and $x \le 2y + 1$. If the first condition fails, then $x \le 2y$, and so $y \ge 3$, since $x \ge 3$ and $y$ is odd. So the second condition holds.

In every case the chosen wall $H'$ deletes $m_v$ and satisfies (a), (b), (c) of (6) with $c = 2$. This proves (10).

Now let $k \ge 2$. Let $(p, q)$ be the midpoint of $m_v$, and let $\delta$ be the least of the numbers $\lfloor q \rfloor$, $h - \lceil q \rceil$, $\lfloor p/2 \rfloor$ and $\lfloor (2h + 1 - p)/2 \rfloor$. These are the distances of $v$ from the bottom, top, left and right sides of $H$, counted in bricks. Every square at $m_v$ has a position $G = (X, Y)$ with $\sigma_h(G) \ge \delta$. For $\lceil q \rceil - 1 \le Y \le \lfloor q \rfloor$ gives $Y + 1 \ge \lceil q \rceil \ge \lfloor q \rfloor$ and $h - Y \ge h - \lfloor q \rfloor \ge h - \lceil q \rceil$, and $|X - p| \le 1$ gives $X + 1 \ge p$ and $2h + 2 - X \ge 2h + 1 - p$. Hence every region of $H$ at $m_v$ has depth at least $\delta$.

\hspace*{12pt}(11) \textit{If $k \ge 2$ and $\delta \le k$, then there are $i, j$ as in (7) such that $T(i, j)$ deletes $m_v$, and every region of $H$ at an element kept by $T(i, j)$ has distance $D_H$ at least $k - 2$ from every region of $H$ at $m_v$.}

We choose $i, j$ by a side of $H$ at distance $\delta$ from $v$, so that $T(i, j)$ leaves a wide margin of $H$ on that side. In the \textit{bottom} case, where $\delta = \lfloor q \rfloor$, let $j = \min(2k - 1, k + 1 + \delta)$ and $i = 2k + (j \bmod 2)$. Otherwise, in the \textit{top} case, where $\delta = h - \lceil q \rceil$, let $j = \max(1, k - 1 - \delta)$ and $i = 2k + (j \bmod 2)$. Otherwise, in the \textit{left} case, where $\delta = \lfloor p/2 \rfloor$, let $i = \min(4k - 1, 2k + 2 + 2\delta)$ and $j = k + ((i + k) \bmod 2)$. Otherwise we are in the \textit{right} case, where $\delta = \lfloor (2h + 1 - p)/2 \rfloor$, and we let $i = \max(1, 2k - 2 - 2\delta)$ and $j = k + ((i + k) \bmod 2)$. These satisfy the conditions of (7), using $k \ge 2$ and $\delta \le k$. In the first two cases $i + j$ is even, $k + 1 \le j \le 2k - 1$ or $1 \le j \le k - 1$, and $i \in \{2k, 2k + 1\}$. In the last two cases $i + j$ is even, $j \in \{k, k + 1\}$, and $2k + 2 \le i \le 4k - 1$ or $1 \le i \le 2k - 2$.

Let $m$ be an element kept by $T(i, j)$, let $s$ be a square at $m$ with position $F = (X, Y)$, and let $s'$ be a square at $m_v$ with position $G = (X', Y')$. By the remark after (2), it is enough to show that the regions $R, R'$ of $H$ containing $s, s'$ satisfy $\text{depth}(R) + \text{depth}(R') \ge k - 2$ and $\|F - G\| \ge k - 2$. The midpoint of $m$ lies in $[i, i + 2h' + 1] \times [j, j + h']$, and so $i - 1 \le X \le i + 2h' + 2$ and $j - 1 \le Y \le j + h'$. Hence the terms of $F$ are at least $j$, $2k - j$, $\lfloor i/2 \rfloor$ and $\lfloor (2h - i - 2h')/2 \rfloor = \lfloor (4k - i)/2 \rfloor$, and so $\text{depth}(R) \ge \min(j, 2k - j, \lfloor i/2 \rfloor, \lfloor (4k - i)/2 \rfloor)$. In each of the four cases this minimum is at least $k - 1 - \delta$. For in the first two cases $\lfloor i/2 \rfloor = k$ and $\lfloor (4k - i)/2 \rfloor \ge k - 1$, and the smaller of $j$ and $2k - j$ is at least $\min(k - 1, 2k - (k + 1 + \delta))$ in the first case and at least $\min(k - 1 - \delta, k + 1)$ in the second. In the last two cases $j \ge k$ and $2k - j \ge k - 1$, and the smaller of $\lfloor i/2 \rfloor$ and $\lfloor (4k - i)/2 \rfloor$ is at least $\min(k - 1, \lfloor (2k - 2 - 2\delta)/2 \rfloor)$. Since $\text{depth}(R') \ge \delta$, we get $\text{depth}(R) + \text{depth}(R') \ge k - 1$.

Now we bound $\|F - G\|$. In the bottom case $Y' \le \lfloor q \rfloor = \delta$, so $\|F - G\| \ge Y - Y' \ge j - 1 - \delta$, which is $k$ if $j = k + 1 + \delta$, and $2k - 2 - \delta \ge k - 2$ if $j = 2k - 1$. In the top case $Y' \ge \lceil q \rceil - 1 = h - \delta - 1$ and $Y \le j + h' = h - 2k + j$, so $\|F - G\| \ge Y' - Y \ge 2k - 1 - \delta - j$, which is $k$ if $j = k - 1 - \delta$, and $2k - 2 - \delta \ge k - 2$ if $j = 1$. In the left case $p < 2\delta + 2$, so $X' \le p + 1$ gives $X' \le 2\delta + 2$, and $\|F - G\| \ge (X - X')/2 \ge (i - 3 - 2\delta)/2$. If $i = 2k + 2 + 2\delta$ this is $k - \frac{1}{2}$, and so $\|F - G\| \ge k$ since it is an integer. If $i = 4k - 1$ this is $2k - 2 - \delta \ge k - 2$. In the right case $p > 2h - 1 - 2\delta$, so $X' \ge p - 1$ gives $X' \ge 2h - 1 - 2\delta$, while $X \le i + 2h' + 2 = i + 2h - 4k + 2$. So $\|F - G\| \ge (X' - X)/2 \ge (4k - 3 - 2\delta - i)/2$, which is $k - \frac{1}{2}$ if $i = 2k - 2 - 2\delta$ and $2k - 2 - \delta$ if $i = 1$. In all cases $\|F - G\| \ge k - 2$.

Finally $T(i, j)$ deletes $m_v$, by the bounds on the midpoints of kept elements. In the bottom case $q < \delta + 1 \le j$, since $2k - 1 \ge k + 1 \ge \delta + 1$. In the top case $q > h - \delta - 1 \ge h - 2k + j = j + h'$, since $j \le 2k - 1 - \delta$. In the left case $p < 2\delta + 2 \le i$, since $4k - 1 \ge 2k + 2 \ge 2\delta + 2$. In the right case $p > 2h - 1 - 2\delta \ge i + 2h' + 1$, since $i \le 4k - 2 - 2\delta$. This proves (11).

For $v$ far from the sides we use a wall with a large hole. Let $k \ge 2$ and $b = k - 1 \ge 1$, so that $h' = h - 2b - 2$, and let $e$ be the even one of $b + 1$ and $b + 2$. Let $Y_0, c_0$ be integers with $1 \le Y_0 \le h'$ and $1 \le c_0 \le 2h' + 1$, and let $Y_1 = Y_0 + b - 1$. For every integer $t$ let $w_t = \max(0, b - |t - Y_1|)$. Thus $w_t = 0$ unless $Y_0 \le t \le Y_0 + 2b - 2$, consecutive values of $w_t$ differ by $1$ for $Y_0 - 1 \le t \le Y_0 + 2b - 1$, and $w_{Y_1} = b$. The wall $L(Y_0, c_0)$ is the subgraph of $H$ defined as follows. Let $R(y) = y$ if $y \le Y_0 - 1$ and $R(y) = y + 2b$ if $y \ge Y_0$. The branch vertex corresponding to $(x, y)$ is the vertex of $H$ in row $R(y)$ and column $x + e$ if $x \le c_0 - 1$, and column $x + e + 2b$ if $x \ge c_0$. The horizontal edges are the paths of $H$ along rows between their ends. The vertical edge at $(x, y)$ is the vertical edge of $H$ at the vertex corresponding to $(x, y)$ if $y \ne Y_0 - 1$. If $y = Y_0 - 1$, the vertical edge at $(x, y)$ is the staircase which has, in each strip $t$ with $Y_0 - 1 \le t \le Y_0 + 2b - 1$, the vertical edge of $H$ in column $x + e - w_t$ if $x \le c_0 - 1$, and in column $x + e + 2b + w_t$ if $x \ge c_0$, together with the horizontal edges of $H$ joining the consecutive ones. So the band of rows $Y_0, \dots, Y_0 + 2b - 1$ of $H$ is not used as rows of $L(Y_0, c_0)$, and the staircases on the left of $c_0$ move one step to the left in each of the first $b$ of these rows, and back to the right in the others. The staircases on the right of $c_0$ do the mirror image. Between them lies a hole shaped like a hexagon, which is the reason for this construction.

\hspace*{12pt}(12) \textit{The wall $L(Y_0, c_0)$ is a wall of height $h'$ of our form, with $c_0$ as given, $y_t = t$ for $t \le Y_0 - 1$, $y_t = Y_0 - 1$ for $Y_0 \le t \le Y_0 + 2b - 1$, $y_t = t - 2b$ for $t \ge Y_0 + 2b$, and $a_t = e - w_t$, $a'_t = e + 2b + w_t$ for every $t$. Let $x_L$ be the largest integer with $x_L \le c_0 - 1$ and $x_L \equiv Y_0 \pmod 2$, and let $B_0 = (x_L + 1 + e + b, Y_1)$. Then all bricks $F$ of $H$ with $\|F - B_0\| \le b$ lie in the same region of $L(Y_0, c_0)$. If moreover $2Y_0 \le h$, then $L(Y_0, c_0)$ satisfies (a), (b), (c) of (6) with $c = 2k$.}

The staircases are paths of $H$. For since $|t + 1 - Y_1| - |t - Y_1| = \pm 1$, consecutive vertical edges of a staircase are in adjacent columns. Moreover $w_t \equiv b + t - Y_1 \equiv t + Y_0 + 1 \pmod 2$, so the column $x + e - w_t$ plus $t$ is odd, since $x \equiv Y_0$ and $e$ is even. The same holds on the right. Since $w_{Y_0 - 1} = w_{Y_0 + 2b - 1} = 0$, each staircase starts and ends at the branch vertices it should join, and it meets rows $Y_0 - 1$ and $Y_0 + 2b$ only in its ends. The columns used lie between $e - b \ge 1$ and $2h' + 1 + e + 3b \le 2h + 1$, and the rows used are at most $h' + 2b = h - 2$. In each row of the band, a staircase uses two adjacent columns, and the staircases on the same side differ by $2$ in $x$. The left staircases use columns at most $c_0 - 1 + e$, and the right ones columns at least $c_0 + e + 2b$. So the staircases are mutually disjoint, and $L(Y_0, c_0)$ has our form with the values given. Here $w_t = 0$ outside the band, and $t - y_t$ is $0$, $t - Y_0 + 1$ or $2b$, which has the parity of $a_t$ and $a'_t$ in each case.

By (4), the move $(\Delta X, \Delta Y)$ of a brick $F = (X, Y)$ satisfies $e - w_Y \le \Delta X \le e + 2b + w_Y$, and $\Delta Y = 0$ if $Y \le Y_0 - 1$, $\Delta Y = Y - Y_0 + 1$ if $Y_0 \le Y \le Y_0 + 2b - 1$, and $\Delta Y = 2b$ if $Y \ge Y_0 + 2b$. In the band, $\Delta Y - b = Y - Y_1$ and $|Y - Y_1| = b - w_Y$. Outside the band $w_Y = 0$ and $|\Delta Y - b| = b$. Hence in all cases
\begin{equation*}
|\Delta X - e - b| + |\Delta Y - b| \le (b + w_Y) + (b - w_Y) = 2b,
\end{equation*}
and $0 \le \Delta Y \le 2b$, with $\Delta Y = 0$ only if $Y \le Y_0 - 1$. Moreover $\Delta X \ge e - w_Y \ge b + 1 - w_Y \ge 1$, and $\Delta X \ge 2$ unless $w_Y = b$, that is, unless $Y = Y_1$.

Let $F = (X, Y)$ be a brick with $\|F - B_0\| \le b$. Then $|Y - Y_1| \le b$, so $Y$ is in the band or $Y = Y_0 - 1$, and $y_Y = Y_0 - 1$. Moreover $|X - (x_L + 1 + e + b)| \le 2b - |Y - Y_1| = b + w_Y$. Hence $X - a_Y = X - e + w_Y \ge x_L + 1 \ge c_0 - 1$ and $X - a'_Y = X - e - 2b - w_Y \le x_L + 1 \le c_0$. In the definition of $P(F)$ in (4), this gives $z = c_0 - 1$ in the first and third cases, and $z = c_0 = x_L + 1$ in the second case. In each case the position of $(z, Y_0 - 1)$ is $(x_L + 1, Y_0 - 1)$, since $x_L + 1 \in \{c_0 - 1, c_0\}$ and $x_L + 1 \equiv Y_0 - 1$. So $P(F) = (x_L + 1, Y_0 - 1)$ for all these bricks, and by (4) they lie in the same region of $L(Y_0, c_0)$.

Now let $2Y_0 \le h$. Any two moves $\mu, \mu'$ satisfy $|\mu_1 - \mu'_1| + |\mu_2 - \mu'_2| \le 4b$ and $|\mu_2 - \mu'_2| \le 2b$, so $\|\mu - \mu'\| \le 2b \le 2c - 2$, and (a) follows as in the proof of (9). Since $h - h' = 2b + 2$, $0 \le \Delta Y \le 2b$ and $1 \le \Delta X \le e + 3b \le 4b + 2$, (5) applies, and the drops are at most $2b$, $2b + 2$, $2b + 1$ and $\lceil (4b + 3)/2 \rceil = 2b + 2$. This gives (b) with $c = 2b + 2 = 2k$. Moreover the top term drops by at most $2b + 1$ unless $\Delta Y = 0$, and the right term drops by at most $2b + 1$ if $\Delta X \ge 2$.

For (c) let $m$ be a deleted element, and let $(x, y)$ be such that $m$ is the vertex $(x, y)$ or an edge at $(x, y)$. Let $s$ be the square $(x, y - 1)$ if $m$ is a vertex or a horizontal edge with $y = Y_1$, and the square $(x, y)$ otherwise. In both cases $s$ is at $m$. Let $F = (X, Y)$ be its position, and suppose that $F$ is a brick of height $h$. We show that each term of $P(F)$ plus $2b + 1$ is at least some term of $F$. The bottom and left terms drop by at most $2b + 1$. If $\Delta Y = 0$, then $Y \le Y_0 - 1$, and the top term of $P(F)$ plus $2b + 1$ is $h' - Y + 2b + 1 = h - 1 - Y \ge Y + 1$, since $2Y + 2 \le 2Y_0 \le h$. So it is at least the bottom term of $F$. It remains to consider the right term when $\Delta X = 1$, which happens only if $Y = Y_1$. Then $m$ is the vertical edge at $(x, Y_1)$, by the choice of $s$, and $X = x + 1$ since $x + Y_1$ is odd. Let $x'' = x - e + b$. Then $x'' \equiv Y_0 \pmod 2$, since $x \equiv Y_1 + 1 = Y_0 + b$. If $0 \le x'' \le c_0 - 1$, then $x = x'' + e - w_{Y_1}$ is the column of the vertical edge in strip $Y_1$ of the staircase at $(x'', Y_0 - 1)$, and so $m$ is kept, which is false. If $x'' < 0$, then $X = x'' + e - b + 1 \le e - b \le 2$, and so $\sigma_h(F) \le 1$. If $x'' \ge c_0$, then $X - a_{Y_1} = x'' + 1 \ge c_0 + 1$, so we are in the second or third case of (4). There $\Delta X = e + 3b$ or $\Delta X \ge X - c_0 = x'' + 1 + e - b - c_0 \ge 2$. So $\Delta X \ge 2$ in all cases where $\sigma_h(F) > 1$. This gives (c), and proves (12).

\hspace*{12pt}(13) \textit{If $k \ge 2$ and $\delta \ge k + 1$, there is a wall $H' \subseteq H$ of height $h - 2k$ which satisfies (a), (b), (c) of (6) with $c = 2k$, and a brick $B$ of $H$ at $m_v$, such that all bricks $F$ of $H$ with $\|F - B\| \le k - 1$ lie in the same region of $H'$.}

Let $b = k - 1$ and $e$ be as above. The square $(\lceil p \rceil - 1, \lceil q \rceil - 1)$ is at $m_v$. Let $B = (X_0, Y_1)$ be its position, so $Y_1 = \lceil q \rceil - 1$ and $\lceil p \rceil - 1 \le X_0 \le \lceil p \rceil$. Let $Y_0 = Y_1 - b + 1$ and $c_0 = X_0 - e - b$. Then $c_0 - 1 \equiv Y_0 \pmod 2$, since $X_0 \equiv Y_1$ and $e$ is even, and so $x_L = c_0 - 1$ in (12), and $B = B_0$. Since $\delta \ge b + 2$, the four side distances of $m_v$ give $\lceil q \rceil \ge b + 2$, $\lceil q \rceil \le h - b - 2$, $p \ge 2b + 4$ and $p \le 2h - 2b - 3$. Hence $Y_0 \ge 2$ and $Y_0 = \lceil q \rceil - b \le h - 2b - 2 = h'$. Moreover $c_0 \ge (2b + 3) - (b + 2) - b = 1$, using $X_0 \ge p - 1$ and $e \le b + 2$, and $c_0 \le (2h - 2b - 3) - (b + 1) - b = 2h'$, using $X_0 \le \lceil p \rceil$ and $e \ge b + 1$. If $2Y_0 \le h$, then (12) shows that $H' = L(Y_0, c_0)$ and $B$ are as required. Otherwise $2(\lceil q \rceil - b) > h$. Now $\rho(m_v)$ has midpoint $(2h + 1 - p, h - q)$, its four side distances are those of $m_v$ in another order, and so the same construction applies to it, with $Y_0$ replaced by $\lceil h - q \rceil - b = h - \lfloor q \rfloor - b$. Since $\lfloor q \rfloor \ge \lceil q \rceil - 1$, twice this number is less than $2h - h - 2b + 2 - 2b \le h$. So the construction gives a wall $L$ for $\rho(m_v)$ with (a), (b), (c), and a brick $B'$ at $\rho(m_v)$ such that all bricks within $\|\cdot\|$-distance $b$ of $B'$ lie in one region of $L$. We take $H' = \rho(L)$, which satisfies (a), (b), (c) by the remark about turned walls, and $B = \rho(B')$. Then $B$ is a brick at $m_v$, since $\rho$ maps the squares at $\rho(m_v)$ to the squares at $m_v$. All bricks within $\|\cdot\|$-distance $b$ of $B$ lie in one region of $\rho(L)$, since $\rho$ preserves $\|\cdot\|$ and maps regions of $L$ to regions of $\rho(L)$. This proves (13).

Now we put the cases together. Let $m_v$ be the element of $H$ that contains $v$. If $r = 0$, then $k = 1$ and (10) gives a wall $H'$ of height $h - 2$ with $d' \ge d - 4$ by (6), which deletes $m_v$. Every vertex $u$ of $H'$ satisfies $u \ne v$, and so $d(u, v) \ge 2 \ge 1$ by (1). So let $r \ge 1$, that is, $k \ge 2$, and let $\delta$ be as above. If $\delta \le k$, let $H' = T(i, j)$ for $i, j$ as in (11). Then $d' \ge d - 4k$ by (7), and every vertex $u$ of $H'$ satisfies $u \ne v$ and lies in an element kept by $H'$. So by (11), every region at $u$ has distance $D_H$ at least $k - 2$ from every region at $v$, and $d(u, v) \ge k$. If $\delta \ge k + 1$, let $H'$ and $B$ be as in (13), and let $O$ be the region of $H'$ that contains the bricks $F$ with $\|F - B\| \le k - 1$. Then $d' \ge d - 4k$ by (6). Let $u$ be a vertex of $H'$. It is at two different regions of $H'$, so at a region $E' \ne O$ of $H'$. As in the proof of (6), some region $E$ of $H$ at $u$ lies in $E'$. So $E$ is not a brick with $\|E - B\| \le k - 1$. Since $B$ is at $m_v$, $\sigma_h(B) \ge \delta \ge k + 1$. If $E$ is the infinite region, then $D_H(E, B) = \sigma_h(B) \ge k + 1$. Otherwise, by (2), $D_H(E, B) = \min(\|E - B\|, \sigma_h(E) + \sigma_h(B)) \ge k$. In both cases $D_H(E, B) \ge k$. Every region of $H$ at $u$ touches $E$ or equals it, and every region of $H$ at $v$ touches $B$ or equals it, since $B$ is at $v$. So every region at $u$ has distance $D_H$ at least $k - 2$ from every region at $v$. Moreover $u \ne v$, since otherwise $E$ would be at $v$ and $D_H(E, B) \le 1$. Hence $d(u, v) \ge k$. In each case $H'$ is a wall of height $h - 2k = h - 2r - 2$ in $H$ with $d' \ge d - 4k = d - 4r - 4$, and no vertex $u$ of $H'$ has $d(u, v) \le r$. This proves (7.3).}\fixnote{fix:7.3-proof-added}{The paper leaves the proof of (7.3) to the reader. The proof above is the one used in the formalization.}

Let $G$ be a graph. A \textit{wall in} $G$ is a subgraph of $G$ which is a wall. Let $H$ be a subgraph of $G$. An \textit{$H$-path in} $G$ is a path with distinct ends, both in $V(H)$, and with no edge or internal vertex in $H$. We shall need the following lemma.

\stmt
(7.4) \textit{For any $p \ge 2$ there exists $r \ge 0$ such that the following holds. Let $H$ be a wall in a graph $G$. Let $P_1, \dots, P_{\binom{p}{2}}$ be mutually vertex-disjoint $H$-paths in $G$, with their $p(p - 1)$ ends mutually at distance $\ge r$. Then $G$ has a minor isomorphic to $K_p$. Moreover, there is an algorithm, with running time $\le C\varepsilon(G)$, where $C$ depends only on $p$, which, given $G$, $H$, and $P_1, \dots, P_{\binom{p}{2}}$ as above, finds a model of $K_p$ in $G$.}

This follows from the main theorem of [16], or more directly from Theorem (4.4) of [20], taking $\mathcal{T}$ to be the set of all separations $(A, B)$ of $H$ of order $< r$ such that $A$ contains no horizontal path of $H$, taking $\Sigma$ to be a sphere, and taking $\Gamma = H$.

Let $H$ be a wall in $G$. For $l \ge 3$, $m \ge 0$ an $(l, m)$-\textit{star over} $H$ in $G$ is a subgraph $P_1 \cup \dots \cup P_l$ of $G$ such that for some vertex $v$,
